\documentclass{article}
\usepackage{amsmath}



\begin{document}
Durante o estudo de progressões, é comum encontrar expressões que envolvem termos elevados a potências.
Considere a seguinte sequência definida por uma relação de recorrência:
\begin{equation}
    a_{n+1} = 2a_n + 3 \quad \text{, sendo $a = 1$}
\end{equation}    
podemos calcular os primeiros termos passo a passo. Alinhe corretamente as equações a seguir:
\begin{align*}
    a_1 &= 1\\
    a_2 &= 2a_1 + 3 = 5\\
    a_3 &= 2a_2 + 3 = 13\\
    a_4 &= 2a_3 + 3 = 29
\end{align*}
Tente usar o ambiente \textit{align} para deixar as igualdades alinhadas no sinal de igual. Depois represente a forma geral da sequência:
$a_n = 2^n - 3$

Por fim, escreva uma frase explicando o que acontece com o valor de 
$a_n$ conforme  $n$ aumenta, e insira uma fórmula \textit{inline} que mostre o termo geral aumentado em uma unidade, como por exemplo $a_{n+1}$.    

\end{document}
